\chapter{Introduction to GUI Programming with Python (Tkinter)}

\section{Introduction: Real-Life Hook}

Think about programs students use every day:

\begin{itemize}
    \item Calculator apps on phones or computers  
    \item Login screens for school portals  
    \item Simple drawing or paint programs  
\end{itemize}

All these are \textbf{GUI (Graphical User Interface) applications}. Tkinter allows Python students to \textbf{create windows, buttons, labels, and interactive apps} easily.

\section{Concept: What is GUI?}

\begin{definition}{Graphical User Interface (GUI)}
A GUI is a \textbf{visual interface that allows users to interact with programs} using buttons, labels, text boxes, and other graphical elements, instead of typing commands in a terminal.
\end{definition}

\begin{itemize}
    \item GUI makes programs \textbf{user-friendly}  
    \item Tkinter is the \textbf{built-in Python library for GUI}  
    \item Simple, beginner-friendly, and suitable for small projects  
\end{itemize}

\section{Setting Up Tkinter}

Tkinter comes \textbf{pre-installed with Python}, so no extra installation is usually needed:

\begin{verbatim}
import tkinter as tk
\end{verbatim}

\section{Creating a Simple Window}

\begin{verbatim}
import tkinter as tk

# Create main window
root = tk.Tk()
root.title("My First GUI")
root.geometry("300x200")  # width x height

# Run the window
root.mainloop()
\end{verbatim}

- \texttt{Tk()} creates the main window  
- \texttt{mainloop()} keeps the window open  

\section{Adding Labels and Buttons}

\begin{verbatim}
import tkinter as tk

def greet():
    label.config(text="Hello, Tkinter!")

root = tk.Tk()
root.title("Greeting App")

label = tk.Label(root, text="Welcome!", font=("Arial", 14))
label.pack(pady=20)

button = tk.Button(root, text="Click Me", command=greet)
button.pack()

root.mainloop()
\end{verbatim}

- \texttt{Label} displays text  
- \texttt{Button} performs an action when clicked  
- \texttt{command} connects button to a function  

\section{Adding Entry Widgets (Text Input)}

\begin{verbatim}
import tkinter as tk

def show_name():
    name = entry.get()
    label.config(text=f"Hello, {name}!")

root = tk.Tk()
root.title("Name App")

entry = tk.Entry(root)
entry.pack(pady=10)

button = tk.Button(root, text="Submit", command=show_name)
button.pack(pady=5)

label = tk.Label(root, text="")
label.pack(pady=10)

root.mainloop()
\end{verbatim}

- \texttt{Entry} allows \textbf{user text input}  
- \texttt{get()} fetches the input value  

\section{Mini-Project: Simple Calculator GUI}

Concept:

\begin{itemize}
    \item Users can \textbf{input two numbers}  
    \item Buttons for \textbf{addition, subtraction, multiplication, division}  
    \item Result displayed in a label  
\end{itemize}

Example:

\begin{verbatim}
import tkinter as tk

def add():
    result = int(num1.get()) + int(num2.get())
    label.config(text=f"Result: {result}")

def subtract():
    result = int(num1.get()) - int(num2.get())
    label.config(text=f"Result: {result}")

root = tk.Tk()
root.title("Simple Calculator")

tk.Label(root, text="Number 1").pack()
num1 = tk.Entry(root)
num1.pack()

tk.Label(root, text="Number 2").pack()
num2 = tk.Entry(root)
num2.pack()

tk.Button(root, text="Add", command=add).pack()
tk.Button(root, text="Subtract", command=subtract).pack()

label = tk.Label(root, text="Result: ")
label.pack(pady=10)

root.mainloop()
\end{verbatim}

- Students learn \textbf{input fields, buttons, functions, and result display}  

\section{Step-by-Step Diagram: Tkinter App Workflow}

\begin{center}
\begin{tikzpicture}[node distance=1.5cm]
\node (start) [draw, ellipse, fill=green!20] {Start GUI App};
\node (window) [draw, rectangle, fill=yellow!20, below of=start] {Create Main Window};
\node (widgets) [draw, rectangle, fill=orange!20, below of=window] {Add Widgets (Labels, Buttons, Entry)};
\node (functions) [draw, rectangle, fill=yellow!20, below of=widgets] {Connect Widgets to Functions};
\node (event) [draw, rectangle, fill=orange!20, below of=functions] {User Interacts with Widgets};
\node (update) [draw, rectangle, fill=yellow!20, below of=event] {Update GUI Based on Interaction};
\node (end) [draw, ellipse, fill=green!20, below of=update] {End App / Close Window};
\draw[->] (start) -- (window);
\draw[->] (window) -- (widgets);
\draw[->] (widgets) -- (functions);
\draw[->] (functions) -- (event);
\draw[->] (event) -- (update);
\draw[->] (update) -- (end);
\end{tikzpicture}
\end{center}

\section{Key Concepts}

\begin{itemize}
    \item GUI allows programs to be \textbf{interactive and visual}  
    \item Tkinter provides \textbf{widgets}: Label, Button, Entry, etc.  
    \item \texttt{command} links buttons to \textbf{functions}  
    \item \texttt{mainloop()} keeps the app running  
    \item Students can build \textbf{calculators, forms, and small apps}  
\end{itemize}

\section{End-of-Chapter Exercises}

\begin{exercise}{Very Short Questions}
\begin{enumerate}
    \item What is a GUI?  
    \item Name two Tkinter widgets.  
    \item What does \texttt{mainloop()} do?  
\end{enumerate}
\end{exercise}

\begin{exercise}{Short Answer Questions}
\begin{enumerate}
    \item Write a Tkinter program to display a label “Hello World”.  
    \item Explain how a button can trigger a function.  
    \item How do you get input from an Entry widget?  
\end{enumerate}
\end{exercise}

\begin{exercise}{Activity-Based Questions}
\begin{enumerate}
    \item Build a GUI app that asks the user for their favorite subject and displays a message.  
    \item Extend the calculator app to include multiplication and division.  
    \item Create a simple login form with username and password fields and a submit button that displays a welcome message.  
\end{enumerate}
\end{exercise}

\section{Teacher Notes}

\marginnote{
\textbf{Learning Objectives:}  
\begin{itemize}
    \item Students can create basic GUI applications with Tkinter  
    \item Students understand widgets, functions, and user interaction  
    \item Students can build small interactive apps  
\end{itemize}

\textbf{Teaching Tips:}  
\begin{itemize}
    \item Start with simple Label and Button examples  
    \item Demonstrate connecting functions to buttons  
    \item Encourage experimentation with fonts, colors, and layouts  
\end{itemize}

\textbf{Assessment:}  
\begin{itemize}
    \item Observe correct creation of widgets and functional buttons  
    \item Check usability and correctness of GUI apps  
\end{itemize}
}